\section{Quality of Service Levels \& Retained Messages}

\subsection{Quality of Service Strategy}
MQTT defines three quality of service (QoS) levels governing delivery guarantees between publisher and broker, and between broker and subscriber. All messages published by the nodes use QoS~0.

\begin{table}[H]
  \caption{QoS Strategy}
  \centering
  \begin{tabular}{llp{0.5\linewidth}}
    \hline
    \textbf{Message Type} & \textbf{QoS} & \textbf{Rationale} \\
    \hline
    Telemetry bursts & 0 (fire and forget) & ACK overhead is disproportionate to benefit; a missed burst is bounded and tolerable; retransmitting a stale sample is actively harmful to the historian. \\
    Full-rate binary packets & 0 & High rate; the acquisition pipeline must never block on the network. \\
    \hline
    \end{tabular}
\end{table}

QoS 2 (exactly once) is not used as the four-way handshake overhead is rarely justified in practice.

\subsection{Retained Messages}
No messages are retained. Retaining high-frequency samples would deliver a burst of stale data to every new subscriber.

On each connection, the firmware publishes an empty retained message to \path{microgrid/<node_id>/status/heartbeat}. This deletes any retained heartbeat left on the broker by earlier firmware versions, and can be removed once no broker holds one.

\subsection{Node Liveness}
Node liveness is determined on the historian side: the Grafana dashboard reports each node's status from the time since its most recent sample. As every node publishes a burst per channel every 5~s, a node that has not reported within a few intervals is treated as offline. This requires no additional MQTT traffic or firmware logic, and also detects failures upstream of the broker connection, such as a stalled acquisition pipeline, that a broker-mediated mechanism would miss.

\paragraph{Deferred} The original design used a retained heartbeat topic and a Last Will and Testament (LWT) message for broker-mediated liveness detection, together with QoS~1 events, alarms, and configuration updates. These were replaced by the historian-side approach above, and no event, alarm, or configuration topics exist.
